\section{Discussion}
\label{sec:discussion}

\curia{} should be viewed primarily as a proof of concept demonstrating that
synteny-constrained embedding comparison yields interpretable cross-species signal
for ncRNA correspondence at genome scale, rather than as a finished annotation
framework optimized for exhaustive sensitivity or production use. Within this
framing, several key observations emerge.

\subsection{Key findings}
We show that cross-species ncRNA correspondence can be recovered at genome scale
using a combination of syntenic constraints and embedding-based similarity. This
enables identification of candidate corresponding loci even in the absence of
strong sequence conservation. For short ncRNAs, correspondence is typically well
localized when supported by annotation, and embedding-based similarity provides a
consistent empirical measure of match confidence. For long ncRNAs, correspondence
is frequently partial and localized, with signal concentrated in discrete
subregions rather than spanning full transcripts. Together, these results support
a model in which cross-species ncRNA correspondence is heterogeneous and modular,
and can be captured using localized comparison in embedding space.

\subsection{Comparison to existing methods}
Existing approaches provide complementary capabilities for ncRNA detection and
structural comparison but do not directly address genome-scale correspondence.
Covariance models require curated family models, while explicit structure
alignment compares supplied RNA sequences rather than locating corresponding
elements from genome and reference annotation alone. Genome-wide scanning
approaches identify candidate RNA elements but do not establish explicit
cross-species correspondence. Methods specifically targeting
lncRNA homology are most directly comparable: slncky infers lncRNA orthologs from
whole-genome alignments using sequence-level conservation signatures, and lncHOME
groups lncRNAs into families based on conserved genomic position combined with
patterns of RNA-binding-protein motifs
\citep{chenEvolutionaryAnalysisMammals2016,huangComputationalPredictionExperimental2024}.
Both rely on primary-sequence or motif-level signals to define homology. In
contrast, \curia{} does not require strong end-to-end sequence conservation
within the ncRNA itself. It combines whole-genome-alignment-derived syntenic
constraints with embedding-based similarity, so
correspondence can be proposed for loci whose primary sequence has diverged beyond
the reach of alignment- or motif-based comparison, without requiring curated
families or explicit structural alignment.

\subsection{Biological implications}
Beyond the modular correspondence pattern noted above, two observations bear on
interpretation. For long ncRNAs, correspondence frequently concentrates in compact
subregions while the surrounding transcript context shows weak or no detectable
similarity. And for short ncRNAs, although embedding-based similarity is strongly
associated with sequence identity, the relationship is not uniform: at intermediate
sequence identity ($\approx 40$--$60\%$), substantial variability is observed,
indicating that loci with comparable sequence similarity can exhibit markedly
different embedding-based similarity.

The island-based representation provides a practical way to capture this behavior:
recurrent embedding-supported islands define candidate corresponding regions, while
intervening regions appear less constrained. Correspondence patterns also vary
across RNA classes: compact RNAs such as tRNAs and snoRNAs show strong and
consistent signal, whereas long ncRNAs more often exhibit partial or fragmented
correspondence
\citep{hezroniPrinciplesLongNoncoding2015,ulitskyEvolutionRescueUsing2016,mattickLongNoncodingRNAs2023}.
Together, these observations support a model in which recurrent candidate elements
are embedded within broader, weakly constrained transcript contexts, emphasizing
the importance of localized comparison in cross-species ncRNA analysis.

\subsection{Limitations}
This study has several limitations related to data availability, modeling
assumptions, and evaluation. First, the lack of reliable ground truth for
cross-species ncRNA correspondence remains a major challenge; existing annotations
are incomplete and heterogeneous across species, and transcript boundaries---particularly
for lncRNAs---are often ambiguous. As a result, evaluation based on
annotation overlap provides only an approximate and conservative estimate of
performance.

Second, \curia{} depends on the availability and quality of genome alignment
chains. Because the search space is restricted to aligned regions, correspondence
cannot be recovered outside chain coverage, and errors or fragmentation in
alignments propagate directly to downstream predictions. In addition, the
candidate-region classification step is adapted from protein-coding gene
projection and may not fully capture the diversity of ncRNA correspondence
patterns.

Third, the method relies on embedding representations derived from RNA foundation
models. Although RiNALMo \citep{penicRiNALMoGeneralpurposeRNA2025} provides a
practical representation, it may not capture all relevant structural features.
The island detector was trained on known compact structured ncRNAs against genomic
background. A genuine lncRNA core whose embedding pattern lies outside that
positive class may therefore be classified as background and never reach the
matcher. Repository held-out-family analyses show broad generalization across
structured-RNA families but substantially weaker recall for lncRNA, consistent
with this failure mode. In particular, \curia{} is primarily sensitive to localized
signal in embedding space and does not model long-range interactions or global RNA
architecture. Because it assumes approximate collinearity within candidate regions
and scores local similarity, it does not capture rearrangements of structured
elements; loci lacking localized embedding signal are unlikely to yield detectable
correspondence.

Fourth, \curia{} does not establish transcriptional activity or molecular
function; predicted regions represent candidate corresponding elements and require
additional transcriptomic or experimental validation. Annotation-discordant calls
have several competing interpretations---missing or imprecise query annotation,
projection or assembly error, correspondence to a different nearby element, or a
false-positive match---which cannot be resolved from interval overlap alone. We
therefore report annotation agreement rather than treating it as either precision
or recall.

Finally, \curia{} does not explicitly distinguish between shared evolutionary
origin and convergent similarity at the level of embedding space. While synteny
strongly biases predictions toward corresponding loci with shared ancestry,
embedding similarity may also capture convergent structural features or recurrent
local motifs. Accordingly, high-confidence matches should be interpreted as
candidate corresponding elements supported by both synteny and embedding
similarity, rather than definitive evidence of shared evolutionary origin.

\subsection{Embedding-based similarity and interpretability}
Embedding representations used in \curia{} are derived from RNA foundation models
trained on large collections of RNA sequences without explicit structural
supervision
\citep{chenInterpretableRNAFoundation2022,akiyamaInformativeRNABase2022,vahabAdvancingNoncodingRNA2025}.
Prior work has shown that such models encode representations associated with RNA
secondary structure and related functional properties, and can support downstream
structure prediction tasks \citep{shenAccurateRNA3D2024}. In this context,
embedding similarity provides an implicit measure of correspondence that is not
tied directly to primary-sequence identity. Empirically, match scores form a
continuous scale separating high-confidence matches from background signal. This
interpretation is broadly consistent with observations from protein language
models, where representations learned from sequence alone organize according to
structural and functional properties \citep{rivesBiologicalStructureFunction2021}.

\subsection{Future directions}
The most informative next step is biological interpretation of the recovered
cores. Predicted elements can be prioritized for secondary-structure analysis,
candidate interaction surfaces and binding partners, transcriptomic support, and
targeted validation. Comparing the number, order, and divergence of cores across
species may reveal how lncRNAs gain, lose, or rearrange functional modules and may
identify recurrent mechanisms that are obscured at whole-transcript scale.

Methodologically, combining core matches across species in a phylogenetic model
could distinguish conserved ancestry from recurrent similarity and reduce the
impact of a single poor assembly or chain set. Multi-scale representations may
also connect locally detectable cores to longer-range RNA architecture.
